Laser light ($\lambda = \laO$) passes a single slit of width $b = \bsO$. The first dark places on either side of the central maximum are at the angle $\alpha$ with $\sin\alpha = \lambda/b$.

\begin{abcliste}
\abc Calculate the angle $\alpha$.
\abc Behind the slit, the photons' momentum sideways spreads by about $\Delta p_y \approx p\,\sin\alpha$. Calculate $\Delta p_y$ and compare $\Delta y\cdot\Delta p_y$ with Heisenberg's limit $h/(4\pi)$, where $\Delta y \approx b$.
\abc The slit is made half as wide. How does the central maximum change?
\end{abcliste}
